\chapter{Time: shift diagnostics and the delivery clock}\label{ch:time-and-clocks}

Most rows of the main and backup copies do not match (\cref{ch:inventory-keys}). The first explanation to test is
time: perhaps the backup holds the same records stored on other dates (\ident{H1}). A shift in time can mean
different things, with different remedies, so this chapter separates them: a lag of the whole series
(\cref{sec:time-global}), offsets of individual rows (\cref{sec:time-rows}), and re-dating of rows that are otherwise
identical (\cref{sec:time-filter}). It then turns to a question every daily figure of the platform depends on: in
which clock were the data's days drawn (\cref{sec:time-clock,sec:time-flips,sec:time-calendar})?

\section{A lag of the whole series}\label{sec:time-global}

If the backup were the main series moved by $k$ days, then for the daily counts $x_t$ (main) and
$y_t$ (backup) the weekly pattern would be rotated by $k \bmod 7$, and after the weekly pattern is removed the
cross-correlation of the series would peak at lag $k$.

\paragraph{Weekday rotation.} For each three-month period with enough common days, let $\bar x_w$ and $\bar y_w$ be the
standardised mean counts on weekday $w$. The rotation $r^\star = \arg\max_r \sum_w \bar x_{w+r} \bar y_w$ is the
shift that best aligns the weekly profiles. It is zero in every three-month period for transactions
($\est{ch04.rot-zero-transactions}{value}$ of $\est{ch04.rot-quarters-transactions}{value}$), contacts
($\est{ch04.rot-zero-call-center-interactions}{value}$ of $\est{ch04.rot-quarters-call-center-interactions}{value}$),
sends and complaints, and in $\est{ch04.rot-zero-digital-events}{value}$ of
$\est{ch04.rot-quarters-digital-events}{value}$ periods for digital events, whose weekly profile is flatter and
noisier (\cref{tab:ch04-global-lag}).

\paragraph{Cross-correlation of the residuals.} Let $u_t$ and $v_t$ be the daily counts minus their weekday means,
standardised. Their sample cross-correlation at lag $h$ is $\hat\rho_{uv}(h) = n^{-1} \sum_t u_t v_{t+h}$. Under
independence of the series, each $\hat\rho_{uv}(h)$ is approximately normal with variance $1/n$, so the band
$\pm z_{1-\alpha/2}/\sqrt{n}$ holds each lag with probability $1 - \alpha$ \autocite{brockwell1991}. Over the
$\est{ch04.ccf-lags}{value}$ lags examined, about $\est{ch04.ccf-expected-outside}{value}$ are expected outside the
band by chance alone. Excluding complaints, the tables have at most $\est{ch04.ccf-outside-max}{value}$ lags outside a
band between $\est{ch04.ccf-band-min}{value}$ and $\est{ch04.ccf-band-max}{value}$, and the largest correlation in
absolute value is $\est{ch04.ccf-max-abs}{value}$, with no peak standing out at any lag. This corrects a statement of the
exploratory notebook, which reported every lag inside the band: a few lags fall outside, as many as chance predicts.
Complaints are the exception that confirms the method: they are identical in both copies, and their correlation at
lag zero is $\est{ch04.ccf-complaints-lag0}{value}$. A global lag is rejected.

\input{../build/generated/tables/ch04-global-lag}

\section{Offsets of individual rows}\label{sec:time-rows}

For a row present in both copies under the same key, the offset is the difference of its dates, backup minus main.
\Cref{tab:ch04-offsets} summarises it for each table. The offset is not a constant: transaction offsets range from
$\est{ch04.offset-min-transactions}{value}$ to $\est{ch04.offset-max-transactions}{value}$ days and campaign-send
offsets from $\est{ch04.offset-min-campaign-sends}{value}$ to $\est{ch04.offset-max-campaign-sends}{value}$ days. It is
often not a whole number of days either: only $\est{ch04.whole-day-digital-events}{value}\,\%$ of shared events and
$\est{ch04.whole-day-campaign-sends}{value}\,\%$ of shared sends differ by an exact multiple of a day. Yet the order of
the shared rows in time is preserved: the rank correlation of \textcite{spearman1904} between the main and backup
dates is at least $\est{ch04.spearman-min}{value}$ in every table. Sequences walking through the same stream of
keys at slightly different speeds produce exactly this signature.

\input{../build/generated/tables/ch04-offsets}

\paragraph{How the offset drifts.} The monthly median offset $o_m$ of each table was fitted by a constant,
a linear trend $o_m = a + bm$, and a piecewise-constant mean with change points found by binary segmentation
\autocite{scott1974,truong2020}. The models are compared by the Bayesian information criterion of
\textcite{schwarz1978},
\begin{equation}\label{eq:bic}
  \mathrm{BIC} = n \log\bigl(\hat\sigma^2\bigr) + k \log n,
\end{equation}
for $n$ months and $k$ parameters, the lower the better. A difference of several units is strong evidence for the
model with the lower value. The constant model loses to the better of the drifting models by
$\est{ch04.bic-gap-transactions}{value}$ for transactions, $\est{ch04.bic-gap-digital-events}{value}$ for events,
$\est{ch04.bic-gap-call-center-interactions}{value}$ for contacts and $\est{ch04.bic-gap-campaign-sends}{value}$ for
sends; for complaints, whose offset is zero, it wins (\cref{tab:ch04-drift-models}). Transactions drift by
$\est{ch04.slope-transactions}{value}$ days per month ($R^2 = \est{ch04.r2-transactions}{value}$) and sends by
$\est{ch04.slope-campaign-sends}{value}$ days per month ($R^2 = \est{ch04.r2-campaign-sends}{value}$); events and
contacts need change points. A fixed storage delay of a set number of days cannot produce drifting, table-specific
offsets of this kind.

\input{../build/generated/tables/ch04-drift-models}

\section{Re-dating of identical rows: the matched filter}\label{sec:time-filter}

The decisive test of a clean re-dating looks at rows whose business content is identical in both copies, whatever
their keys and dates. A re-dating by $k$ days would concentrate their date differences at $k$; independent draws from
the same finite pools of values would spread them over the whole history. For transactions, whose fingerprint
discriminates (\cref{ch:inventory-keys}), the content-identical pairs number $\est{ch04.mf-pairs}{value}$
(\cref{fig:ch04-filter}). Both populations are present. A flat background spreads over
$\est{ch04.mf-support-days}{value}$ days of difference, and a sharp peak sits a few days after zero: the most common
difference, $\est{ch04.mf-mode}{value}$ days, holds $\est{ch04.mf-mode-share}{value}\,\%$ of pairs, and
$\est{ch04.mf-within-week}{value}\,\%$ of pairs lie within a week of zero, where a uniform spread over the observed
range would put $\est{ch04.mf-uniform-week}{value}\,\%$.

\begin{figure}[htbp]
  \centering
  \includegraphics[width=\textwidth]{fig-04-matched-filter}
  \caption{Date difference, backup minus main, between transactions with identical business content. Left: the
  whole range within a little over a year; right: the days around zero. Source:
  \ident{report/build/sources/ch04\_matched\_filter.csv}, recomputed from the pinned parquet files.}
  \label{fig:ch04-filter}
\end{figure}

The peak identifies a small set of transactions that are the same record re-dated. Among the transactions that share
a key, \emph{clones} agree on amount, currency, type, channel, merchant and fraud score; \emph{collisions} share only
the key. There are $\est{ch04.clones}{value}$ clones, $\est{ch04.clones-of-shared}{value}\,\%$ of the shared keys and
$\est{ch04.clones-of-main}{value}\,\%$ of the main copy. Clones and collisions have the same offset distribution
(mean $\est{ch04.clone-offset-mean}{value}$ against $\est{ch04.collision-offset-mean}{value}$ days, standard deviation
$\est{ch04.clone-offset-sd}{value}$ against $\est{ch04.collision-offset-sd}{value}$), so the offset is a property of
the key's position in the stream, not of the record. The customer of a clone agrees across the copies for
$\est{ch04.clones-same-customer}{value}\,\%$ of clones, and its product for $\est{ch04.clones-same-product}{value}\,\%$:
even the re-dated clones were re-keyed.

\paragraph{Verdict on \ident{H1}.} A global lag is rejected; row-level offsets exist but drift smoothly and differ by
table; a clean re-dating holds only for a core of clones. Re-dating the backup would not make the copies comparable.

\section{In which clock were the days drawn?}\label{sec:time-clock}

The source never states the time zone of its timestamps, yet every daily quantity depends on where a day starts:
weekends, holidays, paydays, daily control charts and the hour features of the fraud models. The data can answer the
question themselves.

\begin{proposition}\label{prop:clock}
Suppose events are generated on days of some clock $C$ with a weekly rhythm, and the time of day of an event is drawn
independently of its weekday. Then hour of day and weekday are independent when time is read in clock $C$. Read in a
clock shifted from $C$ by $s$ hours, with $s$ not a multiple of the length $D$ of a day, the weekday boundary cuts
through the hours: events of the last $\lvert s \rvert \bmod D$ hours of each day of $C$ are attributed to the neighbouring weekday, and independence
fails whenever the weekly rhythm is not flat.
\end{proposition}

The clock scan applies \cref{prop:clock}. For each whole-hour shift $s$, it cross-tabulates weekday and hour and
computes the statistic of \textcite{pearson1900} for independence,
\begin{equation}\label{eq:chi2}
  X^2(s) = \sum_{w, h} \frac{(O_{wh} - E_{wh})^2}{E_{wh}},
  \qquad E_{wh} = \frac{O_{w\cdot}\,O_{\cdot h}}{O_{\cdot\cdot}},
\end{equation}
with $\est{ch04.scan-dof}{value}$ degrees of freedom. The generating clock is the shift that minimises $X^2(s)$, where
the statistic per degree of freedom falls to about one. In $\est{ch04.scan-agree}{value}$ of the
$\est{ch04.scan-cells}{value}$ cells of process and market, the minimising shift equals the delivery offset of that
process, $\proto{clock.offset_hours_transactions}$ hours for transactions, digital events and sends and
$\proto{clock.offset_hours_contacts}$ hours for contacts and complaints (\cref{tab:ch04-clock-scan}). At that shift the
statistic per degree of freedom lies between $\est{ch04.scan-best-min}{value}$ and $\est{ch04.scan-best-max}{value}$
for every process except digital events; on the raw timestamps it is at least $\est{ch04.scan-raw-min}{value}$.

Some processes depart from the picture. Digital events reach their minimum at the delivery offset, but their statistic
per degree of freedom stays between $\est{ch04.scan-events-min}{value}$ and $\est{ch04.scan-events-max}{value}$: their
time of day is not independent of the weekday even in their own clock, so \cref{prop:clock}'s premise fails for them.
Satisfaction surveys are minimised at shifts of $\est{ch04.scan-surveys-shift-min}{value}$ to
$\est{ch04.scan-surveys-shift-max}{value}$ hours in the other direction: a survey follows its contact with a delay, so
its timestamp does not carry the delivery clock.

\input{../build/generated/tables/ch04-clock-scan}

\paragraph{The delivery day.} A second, independent test confirms the clock. Every row carries the date of the daily
batch it was delivered in. Read on the raw timestamp, the timestamp's date equals the batch date for as few as
$\est{ch04.clock-raw-min}{value}\,\%$ of rows of a process; shifted by the process's delivery offset, for at least
$\est{ch04.clock-delivery-min}{value}\,\%$ of rows of every process in both copies. For contacts and complaints the
transactions' offset is not enough: shifted by it, only $\est{ch04.clock-minus6-call-center-interactions}{value}\,\%$
of contacts and $\est{ch04.clock-minus6-complaints}{value}\,\%$ of complaints match, which is why the integrity rules
that applied the transactions' offset to these processes reported violations that are artefacts of the clock.

\section{How much the wrong clock moves the calendar}\label{sec:time-flips}

The country series had placed each customer's day on the market's legal time, $\proto{clock.legal_offset_hours_mexico}$,
$\proto{clock.legal_offset_hours_colombia}$ and $\proto{clock.legal_offset_hours_argentina}$ hours from UTC for Mexico,
Colombia and Argentina. The business-date flip rate measures the consequence of that choice: the share of rows whose
date on the legal clock differs from their date on the delivery clock (\cref{tab:ch04-flips}). It is a census share.
Mexico's legal clock coincides with the transactions' delivery clock, and $\est{ch04.flip-mx-flipped}{value}$
Mexican transactions, events or sends change date. For Colombian transactions the share is
$\est{ch04.flip-tx-co}{value}\,\%$ and for Argentine transactions $\est{ch04.flip-tx-ar}{value}\,\%$; across
Argentine processes it ranges from $\est{ch04.flip-ar-min}{value}\,\%$ to $\est{ch04.flip-ar-max}{value}\,\%$, the
highest values belonging to contacts and complaints, whose delivery clock is furthest from the legal one.

\input{../build/generated/tables/ch04-flips}

\section{What the calendar explains}\label{sec:time-calendar}

The country series regressed the logarithm of daily transaction counts on weekday, holiday, long weekend, payday
window, month end, bonus month and a trend, with country fixed effects for the whole bank. On a log scale a
coefficient $\beta$ is a relative effect $e^\beta - 1$, and the intervals use the heteroscedasticity-consistent \HCthree{}
covariance \autocite{mackinnon1985}. They ignore the serial dependence of daily counts (\cref{sec:em-series}), so they
are marked as computed under independence. The regression's p-values form a family fixed by the protocol
($\proto{families.country-calendar-main.size}$ tests on the main copy), and false-discovery control is applied over
the whole family (\cref{sec:em-multiplicity}).

Only $\est{ch04.cal-rejected}{value}$ of the family's tests survive at a false-discovery rate of
$\proto{gates.alpha}$: the weekend terms ($\est{ch04.cal-rejected-weekday}{value}$ of them, Saturday and Sunday in each
market) and the country levels. For the whole bank, Saturday volume differs from Monday's by
$\est{ch04.cal-saturday}{estimate}\,\%$, interval $\est{ch04.cal-saturday}{ci}$, with $q = \est{ch04.cal-saturday}{q}$.
Holidays move volume by $\est{ch04.cal-holiday}{estimate}\,\%$ ($q = \est{ch04.cal-holiday}{q}$). The payday window's
coefficient, $\est{ch04.cal-payday}{estimate}\,\%$, has a p-value of $\est{ch04.cal-payday}{p}$, which would pass a
single test, but its false-discovery adjusted value is $q = \est{ch04.cal-payday}{q}$, and it is not a discovery once
the family is taken into account. The only calendar structure in the data is the weekly rhythm.

On the wrong clock, that rhythm is distorted. In Argentina the legal clock weakens Saturday to
$\est{ch04.sat-ar-local}{value}\,\%$ against $\est{ch04.sat-ar-business}{value}\,\%$ on the delivery clock, and makes
Mondays quieter ($\est{ch04.mon-ar-local}{value}\,\%$ against $\est{ch04.mon-ar-business}{value}\,\%$); Colombia shows
the same distortion, smaller, as its legal clock is closer to the delivery clock. The weekend moves by exactly the hours
the clocks differ, as \cref{prop:clock} predicts.

\section{Implications for the pipelines and the gates}\label{sec:time-implications}

Every daily and hourly quantity of the platform is placed on the delivery clock of its process: the day boundary
of the gold layer, the weekend and hour features of the fraud models, and the hour-grain star on which the readiness
gates are computed (\cref{ch:readiness-gates}). The integrity rules for contacts and complaints must use their own
delivery offset; with the transactions' offset they report violations that do not exist. The backup cannot be
re-dated into the main copy. None of these findings changes a gate's verdict directly; together they define the clock
in which every gate's features and test windows are measured, and an error in that clock would move the weekend and
daypart features the customer and branch gates rely on.

\section{Limitations}\label{sec:time-limits}

The clock is identified by the weekly rhythm; a process without one would leave it undetermined. The scan tests whole
hours only. The cross-correlation band is approximate, and correlations at different lags are not independent, so the
count of lags outside it is a guide, not a test. The flip rate depends on the customer's country, not the transaction's,
because a customer's day is what the calendar describes. The calendar regression's intervals ignore serial dependence,
which makes them too narrow; the conclusion that only the weekend matters rests on false-discovery control of p-values
that share that limitation, and the weekend effects are large enough to survive any plausible widening.
